\section{Historische Einordnung}

Obwohl ausreichend Werke - vor allem Gemälde - existieren, um Italien im \textit{Quattrocento} als Geburtsort der Perspektive zu bestimmen, betont Andersen eine entscheidende Lücke: Es gebe keinerlei Dokumente, die die einzelnen Entwicklungsschritte früherer Generationen belegen. Diese Kluft führte in der Forschung zu zahlreichen Spekulationen. \footnote{Vgl.: \cite[S. 2]{Andersen2007}}
Die Entstehung der Perspektive wird als ein Meilenstein der Menschheitsgeschichte betrachtet. Ich werde mich darauf konzentrieren, einzelne Aspekte zu beleuchten, die zur Geburt der Perspektive beigetragen haben könnten. Sobald wir die historische Perspektive als das notwendige Fundament etabliert haben, lässt sich darauf das mathematisch-historische Gebäude errichten, welches uns schließlich erlaubt, die "`Magie"' dieser Entwicklung -- die \textit{perspektivische Anamorphose} -- fachgerecht zu analysieren und zu begreifen.
Den Ursprung der mathematischen Theorie der Perspektive markiert der italienische Universalgelehrte Leon Battista Alberti, der in seinem Werk \textit{De pictura} um 1435 erstmals ein Modell zum perspektivischen Zeichnen festhielt.\footnote{Vgl.: \cite[S. 1]{Andersen2007}}

Ausgehend von Albertis Konzepten systematisierte der italienische Mathematiker Piero della Francesca im späten 15. Jahrhundert mit seinem Traktat \textit{De prospectiva pingendi} die Grundriss-Aufriss-Methode. \footnote{Vgl.: \cite[S. 14]{Andersen2007}} Mithilfe dieses Verfahrens konstruierte Piero explizit die erste perspektivische Anamorphose - eine gezielte zeichnerische Verzerrung, die erst aus einem speziellen Blickwinkel ihre korrekte Form preisgibt. \footnote{Vgl.: \cite[S. 71-73]{Andersen2007}}

Die neue Darstellungsform wanderte im folgenden Jahrhundert über die Alpen in nördlichere Gebiete ein, darunter auch nach Frankreich und Deutschland.

Angetrieben von der Neugierde an Synthesen aus Geometrie und Optik, eröffnete das  Lehrbuch  \textit{La perspective curieuse, ou magie artificiele des effets merveilleux} (1638) des Minim-Priesters Jean François Niceron (1613–1646) die anamorphotische Kunst auch für Laien: \footnote{Vgl.: \cite[S. 452]{Andersen2007}}

\enquote{TOVTES les parties des Mathematiques ont de rares inuentions, \& des subtilitez qui les ont fait estimer \& cultiuer par les plus beaux esprits de l'antiquité, \& qui les font encore aujourdhuy rechercher par les plus curieux de nostre siecle.} \footnote{\cite[S. 1]{Niceron1663}. dt. Übersetzung: "`Alle Teilbereiche der Mathematik bergen seltene Erfindungen und Feinheiten, derentwegen sie von den feinsten Geistern der Antike geschätzt und gepflegt wurden, und die bewirken, dass sie auch heute noch von den Neugierigsten unseres Jahrhunderts erforscht werden. "'}